% =============================================================
% Distance-Aware Loss for Phonologically-Graded G2P in PT-BR
% IEEE ICASSP 2026 — Anonymous Submission
% =============================================================
%
% COMPILATION:
%   Recommended (best IPA rendering):
%     xelatex main && bibtex main && xelatex main && xelatex main
%
%   Alternative (pdflatex — IPA symbols require a T3-capable font):
%     pdflatex main && bibtex main && pdflatex main && pdflatex main
%
%   If IPA characters show as boxes under pdflatex, uncomment the
%   \usepackage{tipa} block below and replace Unicode IPA with
%   \textipa{} commands.
% =============================================================

\documentclass[conference]{IEEEtran}
\IEEEoverridecommandlockouts

% ---- Encoding & fonts ----------------------------------------
% XeLaTeX path (recommended):
%\usepackage{fontspec}
%\setmainfont{Linux Libertine O}   % or Times New Roman if available

% pdfLaTeX path:
\usepackage[utf8]{inputenc}
\usepackage[T1]{fontenc}

% tipa — required: \textipa{} commands are used throughout
\usepackage{tipa}

% ---- Math ----------------------------------------------------
\usepackage{amsmath}
\usepackage{amssymb}

% ---- Tables --------------------------------------------------
\usepackage{booktabs}
\usepackage{multirow}

% ---- References ----------------------------------------------
\usepackage{cite}

% ---- Hyperlinks (comment out for camera-ready if required) ---
% \usepackage{hyperref}

% ---- Macros --------------------------------------------------
\newcommand{\dpp}{d_{\mathrm{PP}}}        % PanPhon distance
\newcommand{\LCE}{L_{\mathrm{CE}}}        % CE loss
\newcommand{\lam}{\lambda}                % lambda shorthand

% =============================================================
\begin{document}
% =============================================================

\title{Distance-Aware Loss for Phonologically-Graded\\
       Grapheme-to-Phoneme Conversion in Brazilian Portuguese}

\author{%
  \IEEEauthorblockN{[ANONYMOUS]}
  \IEEEauthorblockA{[ANONYMOUS]}
}

\maketitle

% ---- Abstract -----------------------------------------------
\begin{abstract}
Standard grapheme-to-phoneme (G2P) training treats all phoneme
substitutions equally, regardless of articulatory distance.
We propose Distance-Aware (DA) Loss, which penalizes substitutions
proportionally to PanPhon articulatory distance between predicted and
target phonemes, weighted by the model's prediction confidence.
Applied to Brazilian Portuguese with a BiLSTM encoder-decoder,
DA~Loss systematically redistributes errors toward phonologically
closer targets: Class~D (catastrophic) substitutions decrease
19\,\% relative vs.\ a CE~baseline.
Our system achieves PER~0.48\,\% and WER~4.96\,\% on a stratified
28{,}782-word test set---57$\times$ larger than comparable PT-BR
evaluations---with a Wilson 95\,\% CI of $\pm$0.03\,pp.
Evaluation on 31~out-of-vocabulary words yields 100\,\% accuracy
on genuine novel PT-BR words, consistent with phonological rule
generalization beyond memorization.
\end{abstract}

\begin{IEEEkeywords}
grapheme-to-phoneme, distance-aware loss, Brazilian Portuguese,
BiLSTM, phonological features, articulatory distance
\end{IEEEkeywords}

% =============================================================
\section{Introduction}
% =============================================================

Grapheme-to-phoneme (G2P) conversion is a core component of
text-to-speech synthesis, automatic speech recognition, and
multilingual NLP pipelines.
For Brazilian Portuguese, G2P presents well-documented challenges:
graphemic ambiguity (grapheme ``c'' maps to /k/ in \textit{cama}
but /s/ in \textit{cena}; ``r'' maps to /\textipa{\*r}/ in syllable
onset but /x/ in word-final coda~\cite{barbosa2004brazilian}), vowel
neutralization in unstressed positions
(/e/$\leftrightarrow$/\textipa{E}/ and /o/$\leftrightarrow$/\textipa{O}/
merge), and position-dependent coda realization (/x/ in final
position; /\textipa{G}/ before voiced consonants).
Prior work on PT-BR G2P has addressed these challenges with
decision trees, n-gram models, and neural sequence-to-sequence
architectures~\cite{neto2006brazilian}.

Standard models trained with cross-entropy (CE) treat all phoneme
errors equally: predicting /\textipa{E}/ when the target is /e/---a
near-miss differing in one articulatory feature---incurs the same
gradient penalty as predicting /k/ for /a/, an error spanning eight
articulatory features.
This phonological blindness distorts the training signal: the model
learns \emph{that} it erred but not \emph{how severely}.

We address this with \textbf{Distance-Aware (DA) Loss}, which adds a
training signal proportional to the articulatory distance between
predicted and target phonemes, weighted by the model's prediction
confidence.
DA~Loss asks not only ``did you err?'' but also ``how far
phonologically was your error?''
We apply this to Brazilian Portuguese using a BiLSTM encoder-decoder
with Bahdanau attention~\cite{bahdanau2014neural}---deliberately
chosen to isolate the contribution of the loss function from
architectural novelty.

\textbf{Comparison with LatPhon~\cite{chary2025latphon}.}
LatPhon is a 4-layer multilingual Transformer (7.5M parameters, RoPE)
reporting PER~0.86\,\% (Wilson CI $\pm$0.30\,pp) on
$\approx$500~PT-BR words from the same IPA dictionary used here.
Our system achieves PER~0.48\,\% (CI~$\pm$0.03\,pp) on
28{,}782~words.
These test sets differ in size and phonological sampling; this
comparison is indicative rather than confirmatory.
Within these constraints, the Wilson CIs do not overlap (upper
bound 0.51\,\% vs.\ lower bound 0.56\,\%).

\textbf{Contributions:}
\begin{enumerate}
  \item DA~Loss: a phonologically-graded training objective combining
        PanPhon articulatory distance with prediction confidence
        (Sec.~\ref{sec:daloss})
  \item Large-scale stratified evaluation: 28{,}782-word test set
        with Wilson CI 10$\times$ more precise than prior PT-BR work
        (Sec.~\ref{sec:data})
  \item Error quality taxonomy: Class~A--D classification reveals
        systematic redistribution of catastrophic errors
        (Sec.~\ref{sec:analysis})
\end{enumerate}

% =============================================================
\section{Data and Evaluation Protocol}
\label{sec:data}
% =============================================================

\subsection{Corpus}

The training corpus consists of \textbf{95{,}937 grapheme--IPA pairs}
from a Brazilian Portuguese phonetic dictionary.
The input charset covers a--z (excluding k, w, y, which are absent
from the dictionary) plus Portuguese diacritics.
Words containing k, w, or~y are character-OOV and mapped to
$\langle$UNK$\rangle$.
Prior to training, 10{,}252 entries were corrected for ASCII-g
(U+0067) vs.\ IPA-\textipa{g} (U+0261) conflation, required for
correct PanPhon feature lookup.

\subsection{Stratified Split}

The corpus is divided \textbf{60/10/30} (train/val/test) using
stratified sampling over three phonological variables: stress type
(oxy\-tone/paroxytone/proparoxytone), syllable count bin, and
character length bin, yielding approximately 48~strata.
A purely random split risks underrepresenting rare phonological
patterns---proparoxytones account for fewer than 5\,\% of the
corpus---inflating metrics on high-frequency
patterns~\cite{kohavi1995crossvalidation}.
Split quality: $\chi^2\!=\!0.95$ ($p\!=\!0.678$),
Cram\'{e}r $V\!=\!0.0007$---no statistically significant
distributional difference across subsets.
Per-epoch random reshuffling serves a separate purpose: variance
reduction in stochastic gradient estimation~\cite{bottou2010large}.

\subsection{Evaluation Metrics}

\textbf{PER} (Phoneme Error Rate): Levenshtein distance over phoneme
sequences, normalized by reference length~\cite{bisani2008joint}:
\begin{equation}
  \mathrm{PER} = \frac{\displaystyle\sum_i \mathrm{edit}(\hat{y}_i,\,
                 y_i)}{\displaystyle\sum_i |y_i|} \times 100\%
  \label{eq:per}
\end{equation}

\textbf{WER} (G2P Word Error Rate): fraction of words with any
phoneme error (exact-match; equivalent to String Error Rate in
ASR~\cite{bisani2008joint}).

\textbf{Wilson 95\,\% CI}~\cite{wilson1927probable,brown2001interval}:
used throughout in place of Wald intervals, which underestimate
uncertainty near $\hat{p}\!\to\!0$.
For our test set ($\approx$181K reference phonemes), the Wilson CI on
PER is $\pm$0.03\,pp---10$\times$ more precise than
$\pm$0.30\,pp for $\approx$500-word evaluations.

% =============================================================
\section{Architecture}
% =============================================================

We use a \textbf{BiLSTM encoder-decoder with Bahdanau
attention}~\cite{bahdanau2014neural}, established for supervised
G2P~\cite{rao2015g2p}, and chosen to isolate the loss contribution
from architectural novelty.

\textbf{Encoder}: 2-layer BiLSTM over grapheme embeddings.
Each position~$t$ yields $h_t\!=\![\overrightarrow{h}_t;\,
\overleftarrow{h}_t]$, providing full bidirectional context---critical
for resolving graphemic ambiguity where phoneme identity depends on
surrounding characters.

\textbf{Attention}: at each decoder step~$t$:
\begin{align}
  e_{t,j} &= v^\top \tanh(W_h h_j + W_s s_{t-1}), \nonumber \\
  \alpha_{t,j} &= \mathrm{softmax}(e_{t,j}), \quad
  c_t = \textstyle\sum_j \alpha_{t,j}\, h_j
  \label{eq:attention}
\end{align}

\textbf{Decoder}: 2-layer LSTM with teacher forcing during training;
autoregressive at inference.
Three capacity configurations were evaluated
(Table~\ref{tab:configs}).
The medium configuration (9.7M) is the optimal point for DA~Loss;
at 17.2M, model memorization capacity begins to interfere with
phonological regularization (Sec.~\ref{sec:results}).

\begin{table}[t]
  \caption{Architecture configurations}
  \label{tab:configs}
  \centering\small
  \begin{tabular}{lccc}
    \toprule
    Config  & Embedding & Hidden & Parameters \\
    \midrule
    Small   & 128-D & 256-D & 4.3M \\
    Medium  & 192-D & 384-D & 9.7M \\
    Large   & 256-D & 512-D & 17.2M \\
    \bottomrule
  \end{tabular}
\end{table}

% =============================================================
\section{Distance-Aware Loss}
\label{sec:daloss}
% =============================================================

\subsection{Formulation}

\begin{equation}
  \boxed{L = \LCE + \lam \cdot \dpp(\hat{y},\, y) \cdot p(\hat{y})}
  \label{eq:daloss}
\end{equation}

\begin{itemize}
  \item $\LCE = -\log p_\mathrm{correct}$: standard cross-entropy
  \item $\dpp(\hat{y}, y) \in [0,1]$: normalized Euclidean distance
        in PanPhon's 24-dimensional articulatory feature space
        encoding voicing, nasality, place, and manner of
        articulation~\cite{mortensen2016panphon}.
        Example values: p$\leftrightarrow$b: 0.04 (voicing only);
        s$\leftrightarrow$\textipa{S}: 0.15 (place);
        a$\leftrightarrow$k: 0.90 (vowel vs.\ velar stop)
  \item $p(\hat{y}) \in (0,1]$: softmax probability of the
        \emph{predicted} (argmax) phoneme---independent of
        $p_\mathrm{correct}$
  \item $\lam$: coupling weight ($\lam\!=\!0.20$ throughout)
\end{itemize}

\textbf{Design rationale.}
CE monitors $p_\mathrm{correct}$; DA monitors $p(\hat{y})$.
These are independent: when the model is confidently wrong,
$p_\mathrm{correct}$ is low while $p(\hat{y})$ is high.
The confidence factor scales the articulatory penalty with the
model's certainty---maximally penalizing errors that are both
confident and phonologically distant.

\textbf{Bounded signal.}
DA~Loss is upper-bounded by $\lam\!\times\!1.0\!\times\!1.0=0.20$,
while CE can reach $\approx$16 in early training.
DA is effective primarily in the learning transition zone
(CE~0.3--1.5), where the model is actively resolving phoneme
ambiguities.

\textbf{Novelty scope.}
Distance-weighted losses and phonologically-motivated training
signals have been explored in ASR and TTS.
Our contribution is novel in the G2P context through the specific
combination: (1)~PanPhon articulatory distance as gradient signal,
(2)~prediction confidence as weighting factor, and (3)~controlled
coupling via~$\lam$ preserving CE as primary learning axis.

\subsection{Structural Token Override}

PanPhon assigns zero vectors to non-phonemic tokens: syllable
boundary ``\textbf{.}'' and stress marker ``\textbf{'}'' (U+02C8).
Without correction, $d(., \text{'}) = 0$---DA provides no gradient
signal for their confusion.
We apply a post-normalization override: distance~1.0 between any
structural token and any other symbol.
\emph{The override must be applied after Euclidean normalization};
applying it before yields $d\!\approx\!0.25$, not the intended
maximum.

\subsection{$\lambda$ Sweep}

Table~\ref{tab:lambda} sweeps $\lam$ with fixed architecture (4.3M,
no syllable separators).
The inverted-U curve motivates $\lam\!=\!0.20$: too low leaves CE
undifferentiated; too high competes with CE and induces gradient
instability.

\begin{table}[t]
  \caption{$\lambda$ sweep, 4.3M architecture, no separators}
  \label{tab:lambda}
  \centering\small
  \begin{tabular}{cccc}
    \toprule
    $\lambda$ & PER & WER & Behavior \\
    \midrule
    0.05 & 0.62\,\% & 5.36\,\% & Signal too weak \\
    0.10 & 0.63\,\% & 5.35\,\% & Moderate improvement \\
    \textbf{0.20} & \textbf{0.60\,\%} & \textbf{5.14\,\%} & \textbf{Optimal} \\
    0.50 & 0.65\,\% & 5.57\,\% & Over-penalization \\
    \bottomrule
  \end{tabular}
\end{table}

% =============================================================
\section{Experiments}
\label{sec:results}
% =============================================================

\subsection{Error Quality Classes}

Table~\ref{tab:classes} defines four classes by normalized Hamming
distance~$d_H$ in PanPhon 24-feature space.
Training uses Euclidean distance (penalizes bipolar
$+1\!\leftrightarrow\!-1$ feature flips more strongly than neutral
transitions); evaluation uses~$d_H$ (uniform, interpretable).

\begin{table}[t]
  \caption{Phoneme error quality classes}
  \label{tab:classes}
  \centering\small
  \begin{tabular}{clll}
    \toprule
    Class & $d_H$ & Features & Example \\
    \midrule
    A & $0.000$ & 0 --- exact match & --- \\
    B & $\leq\!0.050$ & 1 --- minimal pair & p$\leftrightarrow$b, e$\leftrightarrow$\textipa{E} \\
    C & $\leq\!0.150$ & 2--3 --- same family & s$\leftrightarrow$\textipa{S}, a$\leftrightarrow$\textipa{@} \\
    D & $>\!0.150$  & 4+ --- cross-class & n$\leftrightarrow$\textipa{J}, vowel$\leftrightarrow$stop \\
    \bottomrule
  \end{tabular}
\end{table}

\subsection{Main Results}
\label{sec:mainresults}

Table~\ref{tab:results} presents the main experimental results.
Stratified evaluation across 28{,}782 test words yields Wilson CI
$\pm$0.03\,pp---\textbf{57$\times$ larger} than comparable
PT-BR evaluations~\cite{chary2025latphon}.

\begin{table}[t]
  \caption{Main results. Sep.\ = syllable boundary tokens in output.
           Bold = recommended configurations.}
  \label{tab:results}
  \centering\small
  \begin{tabular}{lcccc}
    \toprule
    System & Params & Sep. & PER & WER \\
    \midrule
    CE baseline              & 4.3M  & ---          & 0.66\,\% & 5.65\,\% \\
    DA $\lambda\!=\!0.1$     & 4.3M  & ---          & 0.63\,\% & 5.35\,\% \\
    DA $\lambda\!=\!0.2$     & 4.3M  & ---          & 0.60\,\% & 5.14\,\% \\
    \textbf{DA} $\boldsymbol{\lambda\!=\!0.2}$ & \textbf{9.7M} & --- & \textbf{0.58\,\%} & \textbf{4.96\,\%} \\
    CE + sep.                & 9.7M  & \checkmark   & 0.52\,\% & 5.79\,\% \\
    \textbf{DA+dist}         & \textbf{17.2M} & \checkmark & \textbf{0.48\,\%} & \textbf{5.33\,\%} \\
    \midrule
    LatPhon~\cite{chary2025latphon} & 7.5M & --- & 0.86\,\% & n/a \\
    \bottomrule
  \end{tabular}
\end{table}

\textbf{Separator trade-off.}
Syllable boundary tokens improve PER ($-17$\,\% to $-20$\,\%) at
the cost of WER ($+6$\,\% to $+8$\,\%).
Each misplaced separator counts as a full word error under
exact-match WER---a structural trade-off independent of architecture
or loss.
\textit{Recommendation}: use the 9.7M DA model (no separators)
for WER-sensitive tasks (NLP, lexicon lookup); use the 17.2M
DA+dist model for PER-sensitive tasks (TTS, forced alignment).

\textbf{LatPhon comparison.}
LatPhon~\cite{chary2025latphon} reports PER~0.86\,\%
($\pm$0.30\,pp, $n\!\approx\!500$).
Our 17.2M DA+dist model achieves PER~0.48\,\%
($\pm$0.03\,pp, $n\!=\!28{,}782$).
CIs do not overlap in the reported scenario.
Test sets differ in size, phonological sampling, and language scope;
this comparison is indicative, not confirmatory.

% =============================================================
\section{Analysis}
\label{sec:analysis}
% =============================================================

\subsection{Dominant Error Patterns}

Table~\ref{tab:confusions} shows the top confusions on the
17.2M~DA+dist model.
Over 60\,\% of errors are vowel neutralizations---phonological
ambiguity inherent to PT-BR.
Structural analysis of the e$\leftrightarrow$\textipa{E} pattern
reveals the cause: the global /e/:/\textipa{E}/ corpus ratio is 7.1:1,
but by position the pre-tonic ratio is 24.9:1 (/e/~dominant) while
the tonic ratio is 0.33:1 (/\textipa{E}/~dominant).
The model learns the strong pre-tonic /e/ bias and overgeneralizes
it to tonic syllables---a corpus distribution artifact.
Resolving this class requires explicit prosodic features (stress
position as input) or a dedicated tonic-vowel classifier.

\begin{table}[t]
  \caption{Top phoneme confusions, 17.2M DA+dist, 28{,}782-word
           test set}
  \label{tab:confusions}
  \centering\small
  \begin{tabular}{lrll}
    \toprule
    Substitution & Count & Group & Mechanism \\
    \midrule
    \textipa{E}$\to$e & 255 & Mid-front vowel & Unstressed neutralization \\
    e$\to$\textipa{E} & 197 & Mid-front vowel & Same (reverse) \\
    \textipa{O}$\to$o & 131 & Mid-back vowel  & Neutralization \\
    i$\to$e           & 121 & High/mid front  & Vowel reduction \\
    o$\to$\textipa{O} & 95  & Mid-back vowel  & Same \\
    \bottomrule
  \end{tabular}
\end{table}

\subsection{Error Quality: DA Loss Redistribution}

Table~\ref{tab:classdist} quantifies the redistribution of errors
along the articulatory axis.
DA~Loss reduces Class~D errors \textbf{19\,\% relative}
(4.3M CE $\to$ 9.7M DA), independently of PER improvement.
Substitutions shift toward phonologically closer targets.
In downstream TTS, Class~B errors (single-feature substitutions,
e.g., /e/$\leftrightarrow$/\textipa{E}/) are perceptually less
salient than Class~D (cross-class, e.g., vowel$\leftrightarrow$stop);
formal perceptual validation (MOS/ABX) remains future work.

\begin{table}[t]
  \caption{Error class distribution. $\dagger$~Class~D inflated by
           structural token confusions (Sec.~\ref{sec:daloss}).}
  \label{tab:classdist}
  \centering\small
  \begin{tabular}{lcccc}
    \toprule
    System & PER & Cls~B & Cls~D & D/errors \\
    \midrule
    CE baseline (4.3M)         & 0.66\,\% & 0.39\,\% & 0.54\,\% & 50.9\,\% \\
    DA $\lambda\!=\!0.2$ (9.7M) & 0.58\,\% & 0.36\,\% & 0.44\,\% & 48.4\,\% \\
    DA+dist (17.2M+sep.)       & 0.48\,\% & 0.29\,\% & 0.53\,$\%^\dagger$ & 47.4\,\% \\
    \bottomrule
  \end{tabular}
\end{table}

\subsection{OOV Generalization}

The 17.2M DA+dist model was evaluated on 31~curated OOV words across
6~categories (Table~\ref{tab:oov}).
The 5/5 result on genuine novel PT-BR words---verified absent from
training---includes correct: palatalization (d+i$\to$d\textipa{Z}),
coda reduction (l$\to$w), rhotacism (rr$\to$x), and nasal vowel
mapping (om$\to$\~{o}).
These patterns are consistent with phonological rule generalization
rather than corpus memorization.

\textbf{Defined failure modes}: (1)~geminate consonants from
Italian/English loans, absent from training data;
(2)~English-phonology anglicisms, genuinely OOV phonologically;
(3)~characters k, w, y---hard character-level OOV.

\begin{table}[t]
  \caption{OOV generalization, 17.2M DA+dist. Phon.\ score =
           average articulatory similarity to reference (0--100\,\%).}
  \label{tab:oov}
  \centering\small
  \begin{tabular}{lcc}
    \toprule
    Category & Correct & Phon.\ Score \\
    \midrule
    PT-BR neologisms        & 6/9\,(67\,\%)   & 97\,\% \\
    Geminate consonants     & 1/5\,(20\,\%)   & 81\,\% \\
    Anglicisms (in-vocab)   & 1/5\,(20\,\%)   & 71\,\% \\
    Char-OOV (k/w/y)        & 0/3\,(0\,\%)    & 68\,\% \\
    \textbf{Real PT-BR OOV} & \textbf{5/5\,(100\,\%)} & \textbf{100\,\%} \\
    Controls (in training)  & 4/4\,(100\,\%)  & 100\,\% \\
    \midrule
    \textbf{Total}          & \textbf{17/31\,(55\,\%)} & --- \\
    \bottomrule
  \end{tabular}
\end{table}

% =============================================================
\section{Conclusion}
% =============================================================

We presented Distance-Aware (DA) Loss, a phonologically-graded
training objective for G2P that penalizes substitution errors
proportionally to PanPhon articulatory distance, weighted by
prediction confidence.
Applied to Brazilian Portuguese, DA~Loss reduces catastrophic
(Class~D) substitutions by 19\,\% relative while achieving
PER~0.48\,\% (Wilson CI $\pm$0.03\,pp) and WER~4.96\,\% on
the largest PT-BR G2P evaluation reported---28{,}782 stratified
words.
Genuine OOV generalization (5/5 novel PT-BR words, 100\,\%)
supports phonological rule acquisition beyond corpus memorization.

DA~Loss is applicable to any language with PanPhon coverage;
effectiveness beyond PT-BR remains to be validated experimentally.
Limitations: geminate reduction for loan words absent from training,
English-phonology anglicisms as genuine phonological OOV, and
homograph disambiguation requiring morphosyntactic context beyond
word-level G2P.

% =============================================================
\bibliographystyle{IEEEtran}
\bibliography{icassp_refs}
% =============================================================

\end{document}
